% Pursuing Stacks (cotes 134-2 to 134-8): archive cote 134-5, site page 39 (PDF page 40), typescript page 226.
% typescript. Draft compared: ps4.tex:1897-1956.
% AI-transcribed by Claude Opus 5.5 (claude-opus-5-5, adaptive thinking, effort high), Message Batches API, batch ps2 (msgbatch_01NVwr1Yuo8qJf5rykeGBXPK), 2026-11-12; reviewed against both supplied scan halves; see CHECK_DECISIONS.tsv.
\scanpage{134-5/39}
\mspage{226}

categories of standard simplices and standard cubes respectively), we
do no have any aspheric functor $A\hand{\to} B$ at hand, and presumably we may
well have that \uline{Asph}$(A,B)$ is empty\hand{,} poor it ! We now assume again that
A and B are test categories, hence $AxB$ is a test category, and the
natural idea for comparison of $\mathrm{Hot}_A$ and $\mathrm{Hot}_B$ is to use the diagram

% Hand-drawn arrows are marked blue; typed objects remain plain.
\diagram{\begin{tikzcd}[column sep=small]
 \& AxB \ar[color=blue,dl,"p_1"'] \ar[color=blue,dr,"p_2"] \& \\
 A \& \& B
\end{tikzcd}}\quad ,

where now $p_1$ and $p_2$ are aspheric (because B and A are aspheric). Thus
we get a corresponding diagram on the corresponding modelizers

\diagram{\begin{tikzcd}[column sep=small]
 \& (AxB)\hat{} \& \\
 A\hat{} \ar[color=blue,ur,"\hand{p_1^*}"] \& \& \deleted{\unsure{Bt}}\, B\hat{} \ar[color=blue,ul,"\hand{p_2^*}"']
\end{tikzcd}}\quad ,

giving rise in turn to equivalences between the corresponding localiza-%
tions

(2)\qquad
\diagram{\begin{tikzcd}[column sep=small]
 \& \mathrm{Hot}_{AxB} \& \\
 \mathrm{Hot}_A \ar[color=blue,ur,"\hand{\overline{p_1^*}}","\hand{\approx}"'] \& \& \mathrm{Hot}_B \ar[color=blue,ul,"\hand{\overline{p_2^*}}"',"\hand{\approx}"]
\end{tikzcd}}\quad ,

These equivalences are compatible with the canonical \xout{funct\unsure{s}} equivalences
with Hot itself, and hence they give rise to an equivalence $\mathrm{Hot}_A\hand{\xrightarrow{\approx}}\mathrm{Hot}_B$
which (up to canonical isomorphism) is indeed the canonical one. This
way for comparing $\mathrm{Hot}_A$ and $\mathrm{Hot}_B$ looks a lot more convenient than the
first one, as the functors $p_1^{*}$ and $p_2^{*}$ which serve as intermediaries
have all desirable exactness properties, and their \inserted{very} description is the
simplest imaginable.

We are very close here to an Eilenberg-Zilber situation, which
will arise more specifically when A and B are both \hand{\uline{\textcolor{black}{strict}}} test categories\hand{,}
i.e. in the corresponding modelizers $A\hat{}$\hand{,}$B\hat{}$ products of models do cor-%
respond to products of the corresponding homotopy types. As seen above,
this implies the same for $A\over{\unsure{x}}{x}B$. Thus, if F is in $A\hat{}$, G in $B\hat{}$, the
object $F\hand{\boxtimes} G$ in $(AxB)\hat{}$ is just a \inserted[typed]{(AxB)-}model for the homotopy type described
respectively by F (in terms of A) and G (in terms of B). As a matter
of fact, the relation
\[ (AxB)_{F\hand{\boxtimes} G} \hand{\simeq} A_{/F}\ x\ B_{/G}\ \hand{,} \]
already noticed before\hand{,} implies that this interpretation holds, indepen-%
dently of strictness. In the classical statement of Eilenberg-Zilber
(as I recall it), we got $A=B$ (both categories A and B being just \hand{$\Delta$});

\margin{39}
